
\subsection{The Adjoint (Reverse‐Mode) Method} \label{sec:adjoint}

The original adjoint calculation method goes back to \cite{cea1986conception} with the use of the Lagrangian of the optimization
problem to compute the derivative of a functional with respect to a shape parameter.


For many-parameter optimization one usually seeks the gradient of a scalar
objective $\Objective(\bm{u},\AutoParameters{})$. 
Using the Lagrangian
%
\[
  \mathcal{L}(\bm{u},\AutoParameters{},\bm{\eta})
  \;=\;
  J(\bm{u},\AutoParameters{})
  +\bigl\langle\bm{\eta},\bm{g}(\bm{u},\AutoParameters{})\bigr\rangle,
\]
%
stationarity with respect to $\bm{u}$ at the solution
$\bm{u}=\bm{u}(\AutoParameters{})$ gives the \emph{adjoint equation}
%
\begin{equation}\label{eq:adjoint}
  \bigl(\partial_{\bm{u}}\bm{g}\bigr)^{\top}\bm{\eta}
  =
  \bigl(\partial_{\bm{u}} \Objective \bigr)^{\top} .
\end{equation}
%
Once $\bm{\eta}$ is known, the gradient
$\nabla_{\!\AutoParameters{}} \Objective$ follows from the \emph{vector–Jacobian product}
%
\begin{equation}\label{eq:vjp}
  \nabla_{\!\AutoParameters{}} \Objective
  \;=\;
  -\,\bigl\langle\bm{\eta},\partial_{\AutoParameters{}}\bm{g}\bigr\rangle
  +\partial_{\AutoParameters{}} \Objective
\end{equation}
or, equivalently,
\begin{equation}\label{eq:grad}
\nabla_{\!\AutoParameters{}} \Objective
\;=\;
-\bm{\AdjointState}^{\mathsf{T}}
\frac{\partial\bm{g}}{\partial\AutoParameters{}}
\;+\;
\frac{\partial \Objective}{\partial\AutoParameters{}} .
\end{equation}


The operation
$\bm{\eta}^{\mathsf{T}}\!\partial\bm{h}(\AutoParameters{})$ in~\eqref{eq:vjp} is the
\emph{vector–Jacobian product} (VJP), the fundamental linear map produced by
\emph{reverse-mode automatic differentiation}.
Reverse mode traverses the computational graph \emph{backwards}, reusing the
solution already computed for the primal Newton solve.

\paragraph{Complexity.}
The cost of solving~\eqref{eq:adjoint} is comparable to one primal
linearization and is \emph{independent of $n_p$}.  
The entire gradient is available after a single VJP.  
The trade-off is memory: all primal
intermediates must be stored, so memory
scales with the depth of the solve.

\paragraph{Automatic–differentiation view.}

\eqref{eq:adjoint}–\eqref{eq:grad} constitute reverse–mode
automatic differentiation: one first executes the \emph{primal} solve to
obtain $\bm{u}$, then runs a reverse sweep that back–propagates the covector
$\bm{\AdjointState}$ through every operation (row–wise evaluation of
$\partial\bm{h}^\mathsf{T}$).  
The cost scales with the
\emph{number of outputs} $n_g$ (often one),
while memory grows with the depth of the computation because intermediate
states must be available for the reverse sweep \cite{MITAD2023}.


Adjoint state techniques allow the use of integration by parts, resulting in a form
which explicitly contains the physically interesting quantity. 
An adjoint state equation is introduced, including a new unknown variable.

The adjoint method formulates the gradient of a function towards its
parameters in a constraint optimization form. 
By using the dual form of
this constraint optimization problem, it can be used to calculate the
gradient very fast. 
A nice property is that the number of computations
is independent of the number of parameters for which you want the gradient. 



% \subsubsection{General case}\label{general_case}


% For a state variable \(\TangentState \in \mathcal{U}\), an optimization variable
% \(v\in\mathcal{V}\), an objective functional
% \(J:\mathcal{U}\times\mathcal{V}\to\mathbb{R}\) is defined. 
% The state
% variable \(u\) is often implicitly dependent on \(v\) through the
% (direct) state equation \(D_{\AutoParameters{}}(u)=0\) (usually the weak form of a partial differential equation),
% thus the considered objective is \(j(v) = J(u_{\AutoParameters{}},v)\), where \(u_{\AutoParameters{}}\) is
% the solution of the state equation given the optimization variables
% \(v\). Usually, one would be interested in calculating \(\nabla j(v)\)
% using the chain rule:

% \[\nabla j(v) = \nabla_{\AutoParameters{}} \Objective(u_{\AutoParameters{}},v) + \nabla_u \Objective(u_{\AutoParameters{}})\nabla_{\AutoParameters{}} u_{\AutoParameters{}}.\]

% Generally the term \(\nabla_{\AutoParameters{}} u_{\AutoParameters{}}\) is often very hard to
% differentiate analytically since the dependence is defined through an
% implicit equation. 
% The Lagrangian functional can be used as a workaround
% for this issue. Since the state equation can be considered as a
% constraint in the minimization of \(j\), the problem

% \[
% \text{minimize }\ j(v) = J(u_{\AutoParameters{}},v)
% \quad
% \text{ subject to }\quad
% \Residual_{\AutoParameters{}}(\TangentState_{\AutoParameters{}}) = 0
% \]

% has an associate Lagrangian functional
% \(\mathcal{L}:\mathcal{U}\times\mathcal{V}\times\mathcal{U} \to \mathbb{R}\)
% defined by

% \[\mathcal{L}(u,v,\AdjointState) = J(u,v) + \langle D_{\AutoParameters{}}(u),\AdjointState\rangle,\]

% where \(\AdjointState \in\mathcal{U}\) is a Lagrange multiplier or \textbf{adjoint state variable} and \(\langle\cdot,\cdot\rangle\) is an inner product
% on \(\mathcal{U}\). 
% The method of Lagrange multipliers states that a
% solution to the problem has to be a stationary
% point of the Lagrangian, namely

% \[\begin{cases}
% d_u\mathcal{L}(u,v,\AdjointState;\delta_u) = d_uJ(u,v;\delta_u) + \langle \delta_u,D_{\AutoParameters{}}^*(\AdjointState)\rangle = 0 & \forall \delta_u \in \mathcal{U},\\
% d_{\AutoParameters{}}\mathcal{L}(u,v,\AdjointState;\delta_{\AutoParameters{}}) = d_{\AutoParameters{}}J(u,v;\delta_{\AutoParameters{}}) + \langle d_{\AutoParameters{}}D_{\AutoParameters{}}(u;\delta_{\AutoParameters{}}),\AdjointState\rangle = 0 & \forall \delta_{\AutoParameters{}}\in\mathcal{V},\\
% d_\AdjointState\mathcal{L}(u,v,\AdjointState;\delta_\AdjointState) = \langle D_{\AutoParameters{}}(u), \delta_\AdjointState\rangle = 0 \quad & \forall \delta_\AdjointState \in \mathcal{U},
% \end{cases}\] where \(d_xF(x;\delta_x)\) is the Gateaux derivative of \(F\) with respect to
% \(x\) in the direction \(\delta_x\). 
% The last equation is equivalent to
% \(D_{\AutoParameters{}}(u) = 0\), the state equation, to which the solution is \(u_{\AutoParameters{}}\).
% The first equation is the so-called adjoint state equation,

% \[
% \langle \delta_u,D_{\AutoParameters{}}^*(\AdjointState)\rangle = -d_uJ(u_{\AutoParameters{}},v;\delta_u) \quad \forall \delta_u \in \mathcal{U},
% \]
% %
% because the operator involved is the adjoint operator of \(D_{\AutoParameters{}}\),
% \(D_{\AutoParameters{}}^*\). 
% Resolving this equation yields the adjoint state
% \(\AdjointState_{\AutoParameters{}}\). 
% The gradient of the quantity of interest \(j\) with
% respect to \(v\) is
% \(\langle \nabla j(v),\delta_{\AutoParameters{}}\rangle = d_{\AutoParameters{}} j(v;\delta_{\AutoParameters{}}) = d_{\AutoParameters{}} \mathcal{L}(u_{\AutoParameters{}},v,\AdjointState_{\AutoParameters{}};\delta_{\AutoParameters{}})\)
% (the second equation with \(u=u_{\AutoParameters{}}\) and \(\AdjointState=\AdjointState_{\AutoParameters{}}\)), thus it
% can be easily identified by subsequently resolving the direct and
% adjoint state equations. The process is even simpler when the operator
% \(D_{\AutoParameters{}}\) is self-adjoint or symmetric since
% the direct and adjoint state equations differ only by their right-hand
% side.

% \subsubsection{Example: Linear case}\label{example_linear_case}

% In a real finite dimensional linear
% programming context, the objective function could be
% \(J(u,v) = \langle Au , v \rangle\), for \(v\in\mathbb{R}^n\),
% \(u\in\mathbb{R}^m\) and \(A \in \mathbb{R}^{n\times m}\), and let the
% state equation be \(K_{\AutoParameters{}} u = b\), with \(K_{\AutoParameters{}} \in \mathbb{R}^{m\times m}\)
% and \(b\in \mathbb{R}^m\).

% The Lagrangian function of the problem is
% \(\mathcal{L}(u,v,\AdjointState) = \langle A u,v\rangle + \left< K_{\AutoParameters{}}u - b,\AdjointState\right>\),
% where \(\AdjointState\in\mathbb{R}^m\).

% The derivative of \(\mathcal{L}\) with respect to \(\AdjointState\) yields the
% state equation as shown before, and the state variable is
% \(u_{\AutoParameters{}} = B_{\AutoParameters{}}^{-1}b\). The derivative of \(\mathcal{L}\) with respect to
% \(u\) is equivalent to the adjoint equation, which is, for every
% \(\delta_u \in \mathbb{R}^m\),

% \[d_u[\langle K_{\AutoParameters{}}\cdot- b, \AdjointState\rangle](u;\delta_u)= - \langle A^\top v,\delta u\rangle \iff \langle K_{\AutoParameters{}} \delta_u,\AdjointState\rangle = - \langle A^\top v,\delta u\rangle \iff \langle K_{\AutoParameters{}}^\top \AdjointState +  A^\top v,\delta_u\rangle = 0 \iff K_{\AutoParameters{}}^\top \AdjointState = - A^\top v.\]
% Thus, we can write symbolically \(\AdjointState_{\AutoParameters{}} = K_{\AutoParameters{}}^{-\top}A^\top v\). 
% The gradient would be

% \[
% \langle \nabla j(v),\delta_{\AutoParameters{}}\rangle = \langle Au_{\AutoParameters{}},\delta_{\AutoParameters{}}\rangle + \langle \nabla_{\AutoParameters{}} B_{\AutoParameters{}} : \AdjointState_{\AutoParameters{}}\otimes u_{\AutoParameters{}},\delta_{\AutoParameters{}}\rangle,
% \]
% where \(\nabla_{\AutoParameters{}} B_{\AutoParameters{}} = \frac{\partial B_{ij}}{\partial v_k}\) is a
% third-order tensor, \(\AdjointState_{\AutoParameters{}} \otimes u_{\AutoParameters{}}\)
% is the dyadic product between the direct and
% adjoint states and \(:\) denotes a double tensor contraction. 
% It is assumed that
% \(B_{\AutoParameters{}}\) has a known analytic expression that can be differentiated
% easily.

% \paragraph{Numerical consideration for the self-adjoint
% case}

% If the operator \(B_{\AutoParameters{}}\) was self-adjoint, \(B_{\AutoParameters{}} = B_{\AutoParameters{}}^\top\), the direct
% state equation and the adjoint state equation would have the same
% left-hand side. 
% In the goal of never inverting a matrix, which is a very
% slow process numerically, a LU decomposition
% can be used instead to solve the state equation, in \(O(m^3)\)
% operations for the decomposition and \(O(m^2)\) operations for the
% resolution. 
% That same decomposition can then be used to solve the
% adjoint state equation in only \(O(m^2)\) operations since the matrices
% are the same.

